\documentclass[12pt]{article}
\usepackage[utf8]{inputenc}
\usepackage[T1]{fontenc}
\usepackage{hyperref}
\title{Ugong Language Resources for K-12 Curriculum Development}
\author{Compiled by Claude Sonnet 5 (claude-sonnet-5) for Andr\'as Kornai}
\date{}
\begin{document}
\maketitle

\section{Introduction}
Ugong (also Gong, Lawa, Ugawng; ISO 639-3 ugo) is a severely endangered Tibeto-Burman language of western Thailand (Uthai Thani and Suphanburi provinces). Speaker estimates vary widely: roughly 50 to 150 speakers in some sources, and about 500 people of Ugong ethnicity in another. Speakers are mostly over fifty and children do not learn it at home. No school materials were found.

\section{Primary Educational Resources}
None found. The language is not known to be taught in schools.

\section{Linguistic and Reference Materials}
The search results name several works, but no online full text was found for them. They include: Mayuree Thawornpat, ``Gong: An endangered language of Thailand'' (Mahidol University PhD, 2006); ``Gong phonological characteristics'', \emph{Mon-Khmer Studies} 37 (2007); David Bradley's papers (1989); and Pusit Rujjanavet's Mahidol MA thesis (1986). Overview pages:

\url{https://en.wikipedia.org/wiki/Gong_language}

\url{https://www.endangeredlanguages.com/lang/ugo}

\url{https://docs.verbix.com/EndangeredLanguages/Ugong}

\section{Religious and Community Resources}
A people-group profile exists: \url{https://people-groups.asiaharvest.org/Thailand/Gong.pdf}. Its content was not examined.

\section{Digital and Audio Resources}
None found. Endangered Languages notes Mahidol documentation projects with community participants, but no outputs were located.

\section{Curriculum Development Context}
With very few older speakers and no teaching materials, any curriculum would begin with documentation and community-led revitalisation. Mahidol University researchers are the main documentation contacts mentioned.
\end{document}